\chapter{多电子原子、Hartree--Fock 与角动量耦合}
\label{chap:many-electron-atoms}

旋转对称性已经把单电子矩阵元压缩成约化矩阵元。多电子原子还多出一个不能忽略的条件：交换两个电子不能产生新的可区分状态，费米统计迫使总态反对称，并由此产生交换能和不同的耦合方案。

\section{反对称多电子空间与 Hartree--Fock 一般结构}

对 $N$ 个相同电子，物理 Hilbert 空间不是普通张量积，而是含自旋的一电子空间 $\mathfrak h$ 的反对称外幂
\begin{equation}
\mathcal H_N=\bigwedge^N\mathfrak h,
\qquad
\mathfrak h=L^2(\mathbb R^3)\otimes\mathbb C^2.
\label{eq:am-fermion-hilbert}
\end{equation}
反对称投影算符为
\[
\mathcal A_N
=\frac{1}{N!}\sum_{\pi\in S_N}\operatorname{sgn}(\pi)P_\pi.
\]
\begin{proposition}[反对称化算符确为正交投影]
$\mathcal A_N$ 满足
\[
\mathcal A_N^\dagger=\mathcal A_N,\qquad
\mathcal A_N^2=\mathcal A_N,
\]
而且任意换位 $\tau\in S_N$ 都满足
$P_\tau\mathcal A_N=-\mathcal A_N$。
\end{proposition}

\begin{proof}
置换算符是酉算符，$P_\pi^\dagger=P_{\pi^{-1}}$，同时
$\operatorname{sgn}(\pi^{-1})=\operatorname{sgn}(\pi)$。把求和指标
$\pi$ 换成 $\pi^{-1}$，便得到
$\mathcal A_N^\dagger=\mathcal A_N$。再计算平方：
\[
\mathcal A_N^2
=\frac1{(N!)^2}
\sum_{\pi,\sigma}
\operatorname{sgn}(\pi\sigma)P_{\pi\sigma}.
\]
对每个固定的 $\rho=\pi\sigma$，共有 $N!$ 对 $(\pi,\sigma)$ 产生
同一个 $\rho$，所以
\[
\mathcal A_N^2
=\frac1{N!}\sum_\rho\operatorname{sgn}(\rho)P_\rho
=\mathcal A_N.
\]
若 $\tau$ 是换位，令 $\rho=\tau\pi$，则
\[
P_\tau\mathcal A_N
=\frac1{N!}\sum_\rho
\operatorname{sgn}(\tau^{-1}\rho)P_\rho
=-\mathcal A_N.
\]
所以 $\mathcal A_N$ 的像只含完全反对称态。反过来，若
$\Psi$ 已完全反对称，则
$P_\pi\Psi=\operatorname{sgn}(\pi)\Psi$，代回定义便有
$\mathcal A_N\Psi=\Psi$。因此它的像恰是完全反对称子空间。
\end{proof}

给定正交自旋轨道 $\{\varphi_a(x)\}$，$x=(\mathbf r,\sigma)$，归一化 Slater 行列式为
\begin{equation}
\Phi(x_1,\ldots,x_N)
=\frac{1}{\sqrt{N!}}\det[\varphi_a(x_i)]_{i,a=1}^{N}.
\label{eq:am-slater-determinant}
\end{equation}
行列式前的 $1/\sqrt{N!}$ 可由正交性直接检查。展开两个行列式后，
积分中出现
\[
\prod_{i=1}^N
\langle\varphi_{\pi(i)}|\varphi_{\sigma(i)}\rangle.
\]
轨道正交使它只在 $\pi=\sigma$ 时非零；共有 $N!$ 个这样的项，
恰好抵消归一化因子的平方。

在无限重核近似下，非相对论 Coulomb 哈密顿量为
\begin{equation}
H
=\sum_{i=1}^{N}
\left[
\frac{\mathbf p_i^2}{2m_e}
-\frac{Ze^2}{4\pi\epsilon_0r_i}
\right]
+\sum_{i<j}\frac{e^2}{4\pi\epsilon_0|\mathbf r_i-\mathbf r_j|}.
\label{eq:am-nelectron-hamiltonian}
\end{equation}
这一定义先于“独立电子”或“屏蔽势”近似。限制变分流形为单个 Slater 行列式后，能量泛函为
\begin{equation}
E_{\rm HF}
=\sum_i\langle i|h|i\rangle
+\frac12\sum_{i,j}(J_{ij}-K_{ij}),
\label{eq:am-hf-energy}
\end{equation}
其中
\[
J_{ij}
=\iint\dd x\,\dd x'\,
|\varphi_i(x)|^2
\frac{e^2}{4\pi\epsilon_0|\mathbf r-\mathbf r'|}
|\varphi_j(x')|^2,
\]
\[
K_{ij}
=\iint\dd x\,\dd x'\,
\varphi_i^*(x)\varphi_j^*(x')
\frac{e^2}{4\pi\epsilon_0|\mathbf r-\mathbf r'|}
\varphi_j(x)\varphi_i(x').
\]
$J$ 是直接 Coulomb 项，$K$ 则完全来自费米子反对称性。

\begin{proof}[Hartree--Fock 方程的变分]
把能量泛函\eqnrefs{eq:am-hf-energy} 在轨道正交归一约束下做变分，逐项算出单体、直接、交换三类贡献，再对角化 Lagrange 乘子矩阵得本征形式的 Fock 方程。

在约束 $\langle\varphi_i|\varphi_j\rangle=\delta_{ij}$ 下极小化
\[
\mathcal L
=E_{\rm HF}
-\sum_{i,j}\Lambda_{ij}
(\langle\varphi_i|\varphi_j\rangle-\delta_{ij}).
\]
对 $\varphi_i^*(x)$ 作独立变分即可（对 $\varphi_i$ 变分给出共轭方程，等价）。

单体项 $\langle i|h|i\rangle=\int\varphi_i^*(x)\,h\varphi_i(x)\dd x$ 对 $\varphi_i^*(x)$ 变分直接给出
\begin{equation*}
 \frac{\delta}{\delta\varphi_i^*}\sum_i\langle i|h|i\rangle
 =h\varphi_i.
\end{equation*}

直接能 $\frac12\sum_{i,j}J_{ij}$ 中，$J_{ij}=\iint|\varphi_i(x)|^2V|\varphi_j(x')|^2$，对 $\varphi_i^*(x)$ 求导得
\[
\frac{\delta J_{ij}}{\delta\varphi_i^*(x)}
=\varphi_i(x)\int\dd x'\,V(x,x')\,|\varphi_j(x')|^2,
\]
因 $J_{ij}=J_{ji}$，求和 $\frac12\sum_{i,j}$ 使 $(i,j)$ 每对被数两次，系数 $\frac12\times2=1$，故直接项总贡献为
\[
\hat J_j\varphi_i(x)
=\left[
\int\dd x'\,
\frac{e^2|\varphi_j(x')|^2}
{4\pi\epsilon_0|\mathbf r-\mathbf r'|}
\right]\varphi_i(x).
\]

交换能 $-\frac12\sum_{i,j}K_{ij}$ 中，$K_{ij}$ 含 $\varphi_i^*(x)$ 与 $\varphi_i(x')$ 两处共轭。对 $x$ 处的 $\varphi_i^*(x)$ 变分（另一半经指标对称 $K_{ij}=K_{ji}$ 与系数 $\frac12$ 合并）给出非局域算符
\[
\hat K_j\varphi_i(x)
=\int\dd x'\,
\frac{e^2\varphi_j^*(x')\varphi_i(x')}
{4\pi\epsilon_0|\mathbf r-\mathbf r'|}
\varphi_j(x),
\]
符号为负，且因 $\hat K_j\varphi_j=\hat J_j\varphi_j$（$i=j$ 时直接项与交换项相消），自相互作用被排除。

约束项变分为 $\sum_{j}\Lambda_{ij}\varphi_j$。合并三类贡献，
\[
\left[h+\sum_{j=1}^{N}(\hat J_j-\hat K_j)\right]\varphi_i
=\sum_j\Lambda_{ji}\varphi_j.
\]

$\Lambda$ 为厄米（约束与泛函皆实），故可用酉矩阵 $U$ 对角化 $\Lambda U=U\,\text{diag}(\varepsilon_i)$。令新轨道 $\widetilde\varphi_i=\sum_j\varphi_j U_{ji}$；酉变换不改变 Slater 行列式（$\det[\widetilde\varphi]=\det[\varphi]\det U$，而 $|\det U|=1$），故总波函数不变。在新基下右端化为对角形式，得到正则 Hartree--Fock 方程
\begin{equation}
\hat f\varphi_i=\varepsilon_i\varphi_i,
\qquad
\hat f=h+\sum_j(\hat J_j-\hat K_j).
\label{eq:am-hf-equation}
\end{equation}
其中 Fock 算符 $\hat f$ 通过 $\hat J_j,\hat K_j$ 依赖全部占据轨道，故本征值问题必须自洽迭代求解。

“独立电子轨道”并非先验假设，而是把完整多电子变分限制在单行列式流形后的自洽场方程；遗漏的行列式之外部分即电子关联能，需由更高阶方法补充。
\end{proof}

\section{氦原子：把变分、自洽场与关联能算到一个数}

Hartree--Fock 方程通常需要数值迭代，但两电子氦原子允许先做一个
完全解析的变分计算。它不会给出高精度结果，却能清楚显示屏蔽、
直接 Coulomb 能和关联能各自从哪里出现。

采用原子单位 $\hbar=m_e=e^2/(4\pi\epsilon_0)=1$。让两个电子占据
同一个空间轨道并组成自旋单重态，试探轨道取
\[
\phi_\zeta(\mathbf r)
=\left(\frac{\zeta^3}{\pi}\right)^{1/2}
\ee^{-\zeta r},
\]
其中 $\zeta>0$ 是待定的有效核电荷。归一化来自
$4\pi\int_0^\infty r^2|\phi_\zeta(r)|^2\dd r=1$。

单个电子的动能与核吸引能分别为
\[
\left\langle-\frac12\nabla^2\right\rangle
=\frac{\zeta^2}{2},
\qquad
\left\langle-\frac Zr\right\rangle=-Z\zeta.
\]
两个电子贡献 $\zeta^2-2Z\zeta$。由于它们自旋相反，交换积分为零，
还需计算直接 Coulomb 积分
\[
J(\zeta)=
\iint\frac{
|\phi_\zeta(\mathbf r_1)|^2
|\phi_\zeta(\mathbf r_2)|^2
}{r_{12}}\dd^3r_1\dd^3r_2.
\]
球对称密度只保留 $1/r_{12}$ 多极展开的 $\ell=0$ 项。角积分给出
\[
\int\dd\Omega_1\dd\Omega_2\frac1{r_{12}}
=\frac{(4\pi)^2}{r_>},
\qquad r_>=\max(r_1,r_2).
\]
把径向区域分成 $r_1>r_2$ 与 $r_2>r_1$ 两个相等部分，
\[
J(\zeta)
=32\zeta^6\int_0^\infty\dd r_1\,r_1\ee^{-2\zeta r_1}
\int_0^{r_1}\dd r_2\,r_2^2\ee^{-2\zeta r_2}
=\frac{5\zeta}{8}.
\]
因此总变分能为
\begin{equation}
E(\zeta)=\zeta^2-2Z\zeta+\frac{5\zeta}{8}.
\label{eq:helium-one-parameter-energy}
\end{equation}
驻值条件
\[
\frac{\dd E}{\dd\zeta}
=2\zeta-2Z+\frac58=0
\]
给出
\[
\zeta_*=Z-\frac5{16},\qquad
E_{\min}=-\left(Z-\frac5{16}\right)^2.
\]
对氦原子 $Z=2$，
\[
\zeta_*=\frac{27}{16},\qquad
E_{\min}=-2.84765625\ {\rm Hartree}.
\]
若忽略电子排斥，两电子都会感受完整核电荷 $Z=2$；变分后得到
$\zeta_*<2$，说明另一个电子屏蔽了部分核电荷。这个单参数结果高于
精确非相对论基态能，因为变分原理给出上界。剩余误差不能仅靠重新
选择一个共同轨道消除：真实波函数会让两个电子的瞬时距离相关，
而单个 Slater 行列式只包含交换关联。

\section{两电子、LS/jj 耦合与 Rydberg 极限}

对两个电子，空间轨道 $a,b$ 的对称/反对称组合为
\[
\Psi_\pm(1,2)
=\frac{\phi_a(1)\phi_b(2)\pm\phi_b(1)\phi_a(2)}
{\sqrt{2(1\pm|S_{ab}|^2)}}.
\]
若轨道正交，$S_{ab}=0$。总费米子波函数必须反对称，所以空间对称态与自旋单重态 $S=0$ 配对，空间反对称态与三重态 $S=1$ 配对。对电子间 Coulomb 相互作用取期望，得到
\begin{equation}
\langle V_{12}\rangle_\pm=J_{ab}\pm K_{ab}.
\label{eq:helium-coulomb-integral}
\end{equation}
这就是 He 类原子单重态--三重态分裂最基本的非相对论来源；更高精度还需叠加自旋轨道、自旋--自旋等 Breit--Pauli 修正。

对多电子原子，当自旋轨道耦合相对电子间静电分裂较弱时，先形成 $\mathbf L=\sum_i\mathbf l_i$ 和 $\mathbf S=\sum_i\mathbf s_i$，再令 $\mathbf J=\mathbf L+\mathbf S$，得到 LS 耦合。重原子强自旋轨道区更自然先形成 $\mathbf j_i=\mathbf l_i+\mathbf s_i$，再把各 $\mathbf j_i$ 耦合成总 $J$。中间耦合不是第三种基本代数，而是同一固定 $J$ 子空间中由真实哈密顿量对不同耦合基的混合。

高激发 Rydberg 电子大部分时间位于离子实之外，长程势趋于 Coulomb 形式。短程核区的多电子效应可压缩进分波相移，因此能级可参数化为
\begin{equation}
E_{nl}
=E_{\rm ion}
-\frac{R_Mhc}{(n-\delta_l)^2},
\label{eq:am-quantum-defect}
\end{equation}
其中 $R_M$ 是含约化质量修正的 Rydberg 常数，$\delta_l$ 是量子缺欠。高 $l$ 轨道受离心势屏蔽而较少进入核区，故 $\delta_l\to0$；低 $l$ 对穿透、交换和离子实极化更敏感。量子缺欠因而是短程多电子散射信息在 Coulomb 长程边界条件中的相移参数化。

\section{原子光谱：从理论到实验}
\label{sec:am-spectroscopy-experiment}

前几节建立了角动量耦合、Hartree--Fock 自洽场和量子缺欠的理论框架。本节把这些结构代入具体原子，给出可与光谱实验直接对比的数值和选择定则，使抽象的表示论成为可操作的原子物理工具。

氢原子是最简单的单电子系统，能级由 Bohr 公式给出
\[
 E_n=-\frac{R_\infty hc}{n^2},
 \qquad
 R_\infty=\frac{m_e e^4}{8\varepsilon_0^2 h^3 c}=1.0973731568\times10^7\,\mathrm{m}^{-1}.
\]
含约化质量修正 $R_M=R_\infty\mu/m_e$ 后，跃迁波长
\[
 \frac1\lambda=R_M\left(\frac1{n_1^2}-\frac1{n_2^2}\right).
\]
Balmer 系（$n_1=2$）的前几条谱线为：
\begin{itemize}
 \item $n_2=3$：H$\alpha$ $\lambda=656.28\,\mathrm{nm}$（红色，可见）
 \item $n_2=4$：H$\beta$ $\lambda=486.13\,\mathrm{nm}$（蓝绿）
 \item $n_2=5$：H$\gamma$ $\lambda=434.05\,\mathrm{nm}$（紫色）
 \item $n_2=6$：H$\delta$ $\lambda=410.17\,\mathrm{nm}$（紫色）
\end{itemize}
Lyman 系（$n_1=1$）在紫外，首线 $\lambda=121.57\,\mathrm{nm}$。Rydberg 常数的实验精度达 $10^{-12}$，是基本常数测定中精度最高的之一。

碱金属的最外层电子在离子实 $(Z-1)e$ 的屏蔽 Coulomb 势中运动，能级由量子缺欠公\eqnrefs{eq:am-quantum-defect} 给出。钠的典型量子缺欠值：
\begin{center}
\begin{tabular}{cccc}
\hline
$l$ & 0 (s) & 1 (p) & 2 (d) \\
\hline
$\delta_l$ & 1.347 & 0.854 & 0.015 \\
\hline
\end{tabular}
\end{center}
由此算出 Na $3s\to3p$ 跃迁（D 线）波长：
\[
 \lambda_D=\frac{1}{R_{\rm Na}\left[\frac1{(3-1.347)^2}-\frac1{(3-0.854)^2}\right]}\approx589.0\,\mathrm{nm},
\]
与实验值 $589.0/589.6\,\mathrm{nm}$ 双线（精细结构分裂 $\Delta\lambda=0.6\,\mathrm{nm}$）一致。高 $l$ 的 $\delta_l\approx0$ 说明 d 轨道几乎不穿入核区，接近纯 Coulomb。

氦原子的两电子组态 $1s2s$ 产生单重态 $^1S_0$ 和三重态 $^3S_1$ 两个态，能量差由\eqnrefs{eq:helium-coulomb-integral} 给出：
\[
 \Delta E=E(^1S_0)-E(^3S_1)=2K_{1s,2s}\approx0.80\,\mathrm{eV}.
\]
交换积分 $K_{1s,2s}$ 的非零值来自两个轨道的空间重叠。实验值：$1s2s\,^3S_1$ 在 $19.820\,\mathrm{eV}$，$1s2s\,^1S_0$ 在 $20.616\,\mathrm{eV}$，差 $0.796\,\mathrm{eV}$。三重态态能量更低（Hund 第一规则）：空间反对称使电子概率密度在接近时减小，Coulomb 排斥降低。

氦基态 $1s^2\,^1S_0$ 的 Hartree--Fock 能量为 $-2.8617\,\mathrm{Hartree}$，实验值为 $-2.9037\,\mathrm{Hartree}$，差 $0.042\,\mathrm{Hartree}\approx1.14\,\mathrm{eV}$ 即电子关联能。此关联能虽仅占总能的 $1.4\%$，但在化学键能尺度（$\sim1\,\mathrm{eV}$）上不可忽略——这正是量子化学需要超越 Hartree--Fock（组态相互作用、耦合簇 CCSD(T) 等）的物理原因。

对若干小原子，Hartree--Fock 基态能量与实验值的对比：
\begin{center}
\begin{tabular}{ccccc}
\hline
原子 & $Z$ & $E_{\rm HF}$ (Hartree) & $E_{\rm exp}$ (Hartree) & 关联能 (Hartree) \\
\hline
He & 2 & $-2.8617$ & $-2.9037$ & $0.0420$ \\
Li & 3 & $-7.4327$ & $-7.4780$ & $0.0453$ \\
Be & 4 & $-14.5730$ & $-14.6674$ & $0.0944$ \\
Ne & 10 & $-128.547$ & $-128.862$ & $0.315$ \\
Ar & 18 & $-526.817$ & $-527.499$ & $0.682$ \\
\hline
\end{tabular}
\end{center}
关联能随电子数增长但占比始终在 $1\%$ 以下。Koopmans 定理给出的 HF 轨道能 $\varepsilon_i$ 近似等于电离能：Ne 的 $2p$ 轨道能 $\varepsilon_{2p}=-0.850\,\mathrm{Hartree}$，实验电离能 $I_p=0.792\,\mathrm{Hartree}$，偏差来自关联和弛豫效应。

电偶极跃迁的选择定则 $\Delta l=\pm1$、$\Delta j=0,\pm1$（$0\not\to0$）、$\Delta m=0,\pm1$ 来自 Wigner--Eckart 定理对 $T^{(1)}_q$（一阶不可约张量）的限制。实验验证：
\begin{itemize}
 \item Na D 线 $3s\to3p$：$\Delta l=+1$，允许；$3s\to4s$：$\Delta l=0$，禁戒（实验中极弱，需磁偶极/电四极机制）。
 \item 氢 Lyman $\alpha$ $1s\to2p$：$\Delta l=+1$，允许，$\lambda=121.6\,\mathrm{nm}$；$1s\to2s$：$\Delta l=0$，电偶极禁戒（双光子跃迁 $1s\to2s$ 需吸收两个光子，速率 $\sim10^{-8}\,\mathrm{s}^{-1}$）。
 \item Ca $4s^2\,^1S_0\to4s4p\,^1P_1$：$\Delta L=1$、$\Delta S=0$、$\Delta J=1$，允许，$\lambda=422.7\,\mathrm{nm}$。
\end{itemize}

Land\'e g 因子\eqnrefs{eq:lande-g-factor} 给出原子在弱磁场中能级分裂 $\Delta E=\mu_B g_J M B$。典型实验值：
\begin{center}
\begin{tabular}{cccccc}
\hline
态 & $L$ & $S$ & $J$ & $g_J$ (理论) & $g_J$ (实验) \\
\hline
$^2S_{1/2}$ & 0 & 1/2 & 1/2 & 2 & $2.0023$ \\
$^2P_{1/2}$ & 1 & 1/2 & 1/2 & 2/3 & $0.666$ \\
$^2P_{3/2}$ & 1 & 1/2 & 3/2 & 4/3 & $1.333$ \\
$^1D_2$ & 2 & 0 & 2 & 1 & $1.000$ \\
$^3P_1$ & 1 & 1 & 1 & 3/2 & $1.501$ \\
\hline
\end{tabular}
\end{center}
理论与实验的微小偏差来自电子 $g_s=2.0023$（QED 修正）和原子核效应。Na D 线在 $B=1\,\mathrm{T}$ 磁场中分裂为 6 条（$D_1$ 线 $^2P_{1/2}\to^2S_{1/2}$ 给出 4 条，$D_2$ 线 $^2P_{3/2}\to^2S_{1/2}$ 给出 6 条），间距 $\mu_B B/h\approx14\,\mathrm{GHz}$，可由 Fabry--Perot 干涉仪直接分辨。

轻原子（$Z\lesssim20$）中电子间静电排斥 $\gg$ 自旋轨道耦合，LS 耦合适用。重原子（$Z\gtrsim80$）中自旋轨道耦合 $\sim$ 静电排斥，趋向 jj 耦合。中间区域需对角化混合矩阵。典型例子：
\begin{itemize}
 \item C $2p^2$：LS 耦合，项 $^3P$, $^1D$, $^1S$，能量间距由 Racah 参数 $B,C$ 决定。
 \item Pb $6p^2$：接近 jj 耦合，$6p_{1/2}^2$, $6p_{1/2}6p_{3/2}$, $6p_{3/2}^2$ 间距由自旋轨道分裂主导。
 \item Xe $5p^5$：中间耦合，需对角化 $9\times9$ 矩阵，LS 和 jj 基都不是好基。
\end{itemize}
这一过渡是原子物理中对称性破缺的典型实例：自旋轨道耦合作为微扰从零增长到与静电排斥同阶，不可约表示的标记从 $(L,S)$ 连续过渡到 $(j_1,j_2)$。

\begin{proof}[三个角动量的重耦合]
三个角动量 \(j_1\otimes j_2\otimes j_3\) 有两种自然耦合方案，由 Racah 6j 符号联系。两种基分别为
\begin{align*}
  |((j_1j_2)J_{12},j_3)JM\rangle&=\sum_{m_1m_2m_3}C^{JM}_{J_{12}M_{12},j_3m_3}C^{J_{12}M_{12}}_{j_1m_1,j_2m_2}|j_1m_1j_2m_2j_3m_3\rangle,\\
  |(j_1,(j_2j_3)J_{23})JM\rangle&=\sum_{m_1m_2m_3}C^{JM}_{j_1m_1,J_{23}M_{23}}C^{J_{23}M_{23}}_{j_2m_2,j_3m_3}|j_1m_1j_2m_2j_3m_3\rangle.
\end{align*}
重耦合系数 \(\langle((j_1j_2)J_{12},j_3)J|(j_1,(j_2j_3)J_{23})J\rangle\) 由 Wigner--Eckart 定理不依赖 \(M\)，定义为 Racah 系数
\begin{equation*}
  W(j_1j_2j_3J;J_{12}J_{23})=\frac{\langle((j_1j_2)J_{12},j_3)J|(j_1,(j_2j_3)J_{23})J\rangle}{\sqrt{(2J_{12}+1)(2J_{23}+1)}}.
\end{equation*}

Racah--Wigner 6j 符号
\begin{equation*}
  \begin{Bmatrix}j_1&j_2&J_{12}\\j_3&J&J_{23}\end{Bmatrix}=(-1)^{j_1+j_2+j_3+J}W(j_1j_2j_3J;J_{12}J_{23}).
\end{equation*}
6j 符号在四角行列式的列置换和上下行交换下不变（共 24 个对称性）；非零条件是四组三角形不等式，对应四面体边长条件。

Racah 1942 给出 6j 符号的有限和闭式：
\begin{equation*}
  \begin{Bmatrix}j_1&j_2&j_3\\j_4&j_5&j_6\end{Bmatrix}=\Delta(j_1j_2j_3)\Delta(j_1j_5j_6)\Delta(j_4j_2j_6)\Delta(j_4j_5j_3)\sum_t\frac{(-1)^t(t+1)!}{\prod_k(t-a_k)!\prod_k(a_k-b_k)!},
\end{equation*}
其中 \(\Delta(abc)=\sqrt{\frac{(a+b-c)!(a-b+c)!(-a+b+c)!}{(a+b+c+1)!}}\) 为三角形系数，求和遍历使所有阶乘非负的整数 \(t\)。对等效电子组态 \(l^n\)，谱项能量由 Racah 参数给出，后者是 Slater--Condon 径向积分 \(F^k\) 经 6j 符号的组合（如 \(d^2\) 的 \(E(^3F)=A-8B\) 等）。四个角动量的两种耦合次序由 9j 符号联系；6j 与 9j 符号把多角动量基底变换写成有限维酉变换。

\(k\) 阶不可约张量算符 \(T^{(k)}_q\) 满足 \([J_z,T^{(k)}_q]=qT^{(k)}_q\)、\([J_\pm,T^{(k)}_q]=\sqrt{k(k+1)-q(q\pm1)}T^{(k)}_{q\pm1}\)。Wigner--Eckart 定理将其矩阵元分离为 CG 几何因子与约化矩阵元：
\begin{equation}
  \langle j'm'|T^{(k)}_q|jm\rangle=\langle j\,m\,k\,q|j'm'\rangle\frac{\langle j'||T^{(k)}||j\rangle}{\sqrt{2j'+1}},
  \label{eq:phys-am-wigner-eckart}
\end{equation}
其中 \(\langle j\,m\,k\,q|j'm'\rangle\) 为纯几何 CG 系数，约化矩阵元 \(\langle j'||T^{(k)}||j\rangle\) 含全部物理信息。非零条件即选择定则：\(m'=m+q\) 来自 \(J_z\) 守恒，三角条件 \(|j-k|\le j'\le j+k\) 来自 \(D^{(k)}\otimes D^{(j)}\) 的分解。
\end{proof}
